%%%%%%%%%%  scan idx 456  |  folio 443  %%%%%%%%%%
\(T_i\) sur \(f_i^{(1)}\). Cet
homomorphisme définit une immersion fermée \(X'\longrightarrow\bfP_X^{d-1}\), régulière d'après 1.8 ii).

D'après 1.8 iii), \(S\) est parfait en tant que \(\cO_X[T_1,\ldots,T_d]\)-Module gradué. En particulier, \(S\) est de présentation finie comme Module gradué. Soit donc:
\[
L''\longrightarrow L'\longrightarrow S\longrightarrow0
\]
une présentation finie graduée de \(S\) sur \(\cO_X[T_1,\ldots,T_d]\). On obtient sur \(\bfP_X^{d-1}\) la suite exacte:
\[
M''\longrightarrow M'\longrightarrow\cO_{X'}\longrightarrow0,
\]
où \(M''=\underline{\Proj}(L'')\), \(M'=\underline{\Proj}(L')\), et sont donc libres de type fini sur \(\bfP_X^{d-1}\). On peut donc écrire pour tout \(n\) le diagramme commutatif:
\[
\begin{tikzcd}[column sep=large,row sep=large]
L''_n \arrow[r] \arrow[d,"\sim"'] &
L'_n \arrow[r] \arrow[d,"\sim"'] &
J^n \arrow[r] \arrow[d] &
0\\
p_*(M''(n)) \arrow[r] &
p_*(M'(n)) \arrow[r] &
p_*(\cO_{X'}(n)) \arrow[r] &
0,
\end{tikzcd}
\]
où \(p\) est le morphisme structural: \(\bfP_X^{d-1}\to X\), et \(L''_n\), \(L'_n\) les composantes de degré \(n\) de \(L''\), \(L'\). La ligne supérieure est exacte par hypothèse; les deux premières flèches verticales sont des isomorphismes à cause des résultats classiques pour les images directes dans \(\cO_P(n)\), \(P\) étant un fibré projectif (EGA III, 2.1.15); enfin, la dernière ligne est exacte si \(n\) est assez grand, car le noyau de \(M'\to\cO_{X'}\) est un \(\cO_P\)-Module parfait, et on peut appliquer le lemme rappelé plus haut, qui montre que ce noyau, tensorisé par \(\cO_P(n)\) pour \(n\) assez grand, est \(p_*\)-acyclique; de même pour le noyau de \(M''\to M'\). Il en résulte alors que pour \(n\) assez grand, \(J^n\longrightarrow p_*(\cO_{X'}(n))\) est un isomorphisme. Or \(p_*(\cO_{X'}(n))=f_*(\cO_{X'}(n))\); on a
